\section{Conclusion}
\label{sec:conclusion}

The core challenge of scientific discovery lies in searching for optimal solutions within vast, structured and open-ended design spaces under a limited budget of costly evaluations. Existing paradigms exhibit complementary limitations: large language models (LLMs) possess strong structured generative capabilities and domain priors, but cannot reliably estimate the true value of external objectives or their associated uncertainty. Bayesian optimisation, by contrast, delivers uncertainty-aware experimental decisions from sparse observations, yet struggles to autonomously propose valid candidates in complex discrete spaces. This work presents the Large Discovery Model (LDM), an experiment-grounded recurrent architecture that deeply couples the generative prior of LLMs with the value signal from a Gaussian process surrogate. Framing scientific design as a sequential inverse decision problem under an evaluation budget, LDM operates a closed loop of generation, evaluation and update that simultaneously expands the search frontier and allocates experimental resources precisely.

At the theoretical and algorithmic level, LDM establishes a KL-regularised, acquisition-guided optimisation framework and derives a closed-form acquisition-tilted optimal search distribution. This formulation extends the classical dichotomy between exploitation and exploration into a tripartite scientific search regime of exploitation, exploration and discovery, unifying the generative prior and empirical value signal within a single mathematical framework. We further provide a complete regret decomposition that explicitly delineates the performance boundaries of generative coverage gap, surrogate fitting error and sampling strategy suboptimality. Algorithmically, the framework supports two candidate generation paradigms, direct generation and indirect parameterisation, and is complemented by a decomposed feedback mechanism to guide targeted iterative refinement, making it adaptable to design spaces of varying structure. Building on this core, we introduce a post-training amortisation scheme: through reasoning-augmented supervised fine-tuning, the high-budget test-time search policy is amortised into the model weights, substantially reducing inference-time computational cost while preserving search quality.

Empirical evaluation across three heterogeneous domains, including the auto-research scenario of neural-network training programme, antibody CDRH3 sequence design, and multi-objective molecular optimisation, systematically validates the generality of the LDM framework and its core findings. First, the structured generative capacity of LLMs and the uncertainty calibration of Bayesian surrogates exhibit strong complementarity, with their integrated performance consistently outperforming both LLM-only and BO-only baselines. Second, the discovery mechanism is central to breaking performance plateaus: purely local optimisation converges rapidly to local optima, whereas expansion of the search frontier continuously lifts the attainable performance ceiling. Third, both test-time compute scaling and post-training fine-tuning reliably improve search efficiency, and the optimal generation paradigm correlates with the strength of domain priors. Fine-tuning experiments further demonstrate stable gains across cross-target and cross-task transfer, indicating that the model learns generalisable discovery decision logic rather than task-specific memorisation.

Nevertheless, the current framework is subject to several constraints and limitations. The Gaussian process surrogate incurs cubic computational complexity in the number of observations, limiting its scalability under large experimental budgets. Test-time search relies on batch candidate sampling, which entails substantial LLM inference overhead. The kernel functions and feature representations of the surrogate still require manual customisation per domain, resulting in limited automation for cross-domain transfer. Additionally, framework performance is bounded by the coverage of the LLM’s pretrained domain knowledge; in domains with weak priors, it is difficult to achieve substantial superiority over specialised methods. Finally, the current closed loop is updated solely based on in-process experimental observations, and has not yet systematically incorporated external existing knowledge and data, leaving a large body of unknown knowns unutilised in the search process.

For future work, a core direction is to address the utilisation of unknown knowns. We will introduce memory retrieval mechanisms and federated discovery frameworks to incorporate external knowledge bases, published experimental results and data from parallel exploration processes into the search closed loop. Combined with techniques such as sparse Gaussian processes \citep{snelsonSparseGaussianProcesses2005, titsiasVariationalLearningInducing2009}, this will simultaneously resolve the dual challenges of limited context capacity and high computational complexity as observational data scales. Building on this, we will explore end-to-end surrogate specification, in which the LLM autonomously learns kernel functions and prior settings for the surrogate, eliminating manual adaptation across domains. We will also advance discovery-oriented post-training paradigms, exploring more sophisticated training paths such as preference learning and reinforcement learning to further internalise the acquisition-tilted decision logic into the model’s generative distribution. This would endow the model with native capability to balance exploitation, exploration and discovery, ultimately enabling more efficient open-ended scientific discovery.
